% ============================================
% 在线论文文档自动生成工具 — 软件需求规格说明书（完整版）
% 由全组六名成员共同编写，此处合并为完整文档
% ============================================

\documentclass{../softeng}

\begin{document}

\doccover{软件需求规格说明书}
{在线论文文档自动生成工具}
{V2026.07}

\tableofcontents
\newpage

\begin{revrecord}
V1.0 & 2026-07-15 & 全组 & 全组 & 初稿完成 \
\hline
\end{revrecord}
% ============================================
% 第1章：引言 + 任务概述 + 数据描述
% 负责人：A
% ============================================

\section{引言}

\subsection{编写目的}

本文档为在线论文文档自动生成工具（Thesis Generator）软件需求规格说明书，用于完整定义系统的全部功能性与非功能性需求，明确项目开发边界、业务规则、数据规范和运行约束。

本文档的预期读者包括：
\begin{itemize}
    \item 项目经理：用于制定项目计划和跟踪需求实现情况；
    \item 系统设计人员：作为概要设计和详细设计的技术依据；
    \item 前后端开发工程师：明确各自负责模块的功能边界和接口约定；
    \item 测试人员：作为编写测试用例和验收测试的标准参考。
\end{itemize}

\subsection{项目背景}

本项目为高校大二软件工程实训项目（CST328 项目实训），开发单位为在校实训开发小组，共六名成员。项目目标是开发一套Web端在线论文编辑与标准化文档生成系统，服务于高校学生在课程论文、毕业论文、项目设计类论文撰写过程中的规范化需求。

当前高校学生在撰写论文时面临以下痛点：论文模板格式不统一、排版操作繁琐、参考文献格式手动调整耗时、文档导出格式错乱、课程论文批阅流程缺乏线上化支持。本系统旨在通过Web平台实现从模板分发、在线写作、预览核查、文档导出到教师线上批改的一体化论文生成工具，降低学生论文撰写的格式负担，提高教师批阅效率。

系统独立部署运行，无第三方外部系统强制对接，仅依托Web服务、关系型数据库、Redis缓存和文件存储完成完整业务闭环。

\subsection{术语与定义}

\begin{itemize}
    \item \textbf{论文模板}：预设的论文结构框架，包含封面字段配置、章节结构、排版格式参数（字体、字号、行距、页边距等），学生创建论文时自动套用；
    \item \textbf{模板版本}：同一模板的多次修订记录，每次修改生成新版本，学生新建论文自动使用最新启用版本；
    \item \textbf{富文本编辑器}：基于Web浏览器的所见即所得（WYSIWYG）在线编辑器，支持文字格式化、图片嵌入、表格插入、引用标记等功能；
    \item \textbf{批注}：教师在学生论文指定文字位置添加的修改意见或评论，与论文原文位置绑定；
    \item \textbf{批阅}：教师对学生提交的课程论文进行审查、添加批注、评定分数和等级的操作流程；
    \item \textbf{提交}：学生完成论文编辑后将论文发送至教师端等待批阅的操作，提交后论文锁定不可编辑；
    \item \textbf{退回}：教师将不符合要求的论文返回给学生重新修改的操作，退回后论文解锁恢复编辑；
    \item \textbf{DOCX/PDF导出}：将学生在线编辑的论文内容按照模板格式配置生成标准Word文档或PDF文件的功能。
\end{itemize}

\subsection{参考资料}

\begin{enumerate}
    \item GB/T 9385-2008 计算机软件需求规格说明规范；
    \item 2026软件工程实训项目任务书；
    \item Web前后端分离开发设计规范；
    \item Apache POI 官方文档——Java操作Microsoft Office文档的开源库；
    \item GB/T 7714-2015 信息与文献 参考文献著录规则；
    \item Spring Boot 官方参考文档；
    \item PostgreSQL 数据库设计开发规范。
\end{enumerate}

\section{任务概述}

\subsection{目标}

搭建一套Web端在线论文编辑与标准化文档生成平台，区分系统管理员、学生、教师三类用户角色，分别赋予模板管理、论文编辑导出、课程论文线上批阅等对应权限，实现以下核心业务目标：

\begin{enumerate}
    \item 管理员端：支持按学院维度对论文模板进行分类管理，可自定义发布多类型模板（毕业论文、课程论文、项目论文），模板内置封面、中英文摘要、多级章节正文、参考文献等固定结构，统一约束字体、行距、页边距等排版格式，支持模板版本管理与在线预览；
    \item 学生端：可按需选用模板，依托结构化章节大纲在线完成论文内容编写（含文字排版、图片嵌入、表格编辑、参考文献管理），支持云端草稿存储、全文Web预览、一键导出标准DOCX和PDF格式文档，课程论文可提交至教师端等待批阅；
    \item 教师端：可查看学生提交的课程论文内容，在文档指定位置添加批注修改意见，填写整体评语并评定分数与等级，支持退回论文供学生修改重提，留存历次批阅历史记录。
\end{enumerate}

\subsection{运行环境}

\begin{itemize}
    \item \textbf{服务端操作系统}：Windows Server 2019 / Linux CentOS 7 及以上；
    \item \textbf{Web服务容器}：Spring Boot内嵌Tomcat / Nginx反向代理；
    \item \textbf{数据库}：PostgreSQL 16+，存储全部业务数据；
    \item \textbf{缓存与队列}：Redis 7+，用于Token缓存、热点数据缓存、异步任务队列；
    \item \textbf{文件存储}：服务器本地文件系统，存储上传的论文图片资源及生成的DOCX/PDF文件；
    \item \textbf{客户端环境}：Chrome（最新版）、Edge（最新版）、Firefox（最新版）等主流现代浏览器。
\end{itemize}

\subsection{条件与限制}

\begin{enumerate}
    \item 系统仅开放管理员、学生、教师三类角色，不开放游客的论文编辑和导出权限；
    \item 用户登录采用JWT Token鉴权机制，所有API接口需携带有效Token，按角色严格区分数据访问权限；
    \item 学生论文编辑过程中上传的图片单张不得超过5MB，格式限制为jpg/png/gif；
    \item DOCX和PDF文档生成过程涉及服务端文件IO操作，单次导出操作应在60秒内完成；
    \item 系统为独立部署的Web应用，不依赖第三方外部系统接口，所有功能在封闭环境中完成；
    \item 服务器硬件资源有限（实训环境），需合理控制并发导出任务数量和文件存储空间。
\end{enumerate}

\section{数据描述}

\subsection{静态数据}

系统长期持久存储且不频繁变更的基础数据包括：

\begin{itemize}
    \item \textbf{用户账号信息}：用户名、加密密码、真实姓名、角色标识（ADMIN/STUDENT/TEACHER）、所属学院；
    \item \textbf{学院基础信息}：学院名称、学院编号；
    \item \textbf{论文模板信息}：模板名称、论文类型（毕业论文/课程论文/项目论文）、关联学院、启用状态；
    \item \textbf{模板版本配置}：封面字段配置（JSON）、章节结构定义（JSON）、排版格式参数（JSON）；
    \item \textbf{题目讨论与公告}：系统公告内容、讨论帖子内容（本系统暂不涉及）。
\end{itemize}

\subsection{动态数据}

系统运行过程中实时产生、频繁变更的数据包括：

\begin{itemize}
    \item \textbf{用户登录会话}：JWT Token及其Redis缓存，具有时效性（24小时过期）；
    \item \textbf{论文内容数据}：学生在线编辑的论文章节HTML内容、草稿自动保存版本；
    \item \textbf{参考文献条目}：学生逐条录入的参考文献信息（作者、标题、刊物、年份等）；
    \item \textbf{上传图片记录}：学生在论文中嵌入的图片文件及元数据；
    \item \textbf{提交记录}：学生历次提交论文的时间戳和版本号；
    \item \textbf{批注与批阅记录}：教师在论文指定位置添加的批注意见、整体评语、分数和等级；
    \item \textbf{导出任务状态}：DOCX/PDF生成任务的进度和结果。
\end{itemize}

\subsection{数据库介绍}

采用PostgreSQL关系型数据库存储全部业务数据。数据库名称为\verb|thesis_generator|，设计包含以下核心数据表：

\begin{itemize}
    \item \verb|sys_user|：系统用户表，存储全部用户账号、角色和所属学院信息；
    \item \verb|sys_college|：学院信息表，存储学院基础数据；
    \item \verb|template|：论文模板表，存储模板基本信息和启用状态；
    \item \verb|template_version|：模板版本表，存储每个模板的各版本配置数据（封面字段、章节结构、格式参数），使用JSONB类型字段以支持灵活的配置结构；
    \item \verb|thesis|：论文主表，存储每篇论文的元信息（标题、状态、关联学生和模板版本）；
    \item \verb|thesis_section|：论文章节表，存储每篇论文各章节的结构化信息和HTML内容；
    \item \verb|thesis_reference|：参考文献表，存储每篇论文的参考文献条目；
    \item \verb|thesis_image|：图片资源表，记录学生在论文中上传的图片文件信息；
    \item \verb|submission|：提交记录表，记录学生历次提交论文的时间信息；
    \item \verb|annotation|：批注表，存储教师在论文指定位置添加的批注意见；
    \item \verb|review_record|：批阅记录表，存储教师的评语、分数、等级和操作类型。
\end{itemize}

各数据表之间通过外键建立关联约束：用户关联学院、论文关联用户和模板版本、章节和参考文献关联论文、批注和批阅记录关联论文等。

\subsection{数据词典}

以下是系统核心数据元素的简要词典：

\begin{table}[htbp]
    \centering
    \caption{核心数据词典}
    \begin{tabular}{|c|c|c|c|}
        \hline
        \textbf{数据项} & \textbf{所属表} & \textbf{数据类型} & \textbf{说明} \\
        \hline
        用户ID & sys\_user & BIGINT & 自增主键，唯一标识用户 \\
        \hline
        用户名 & sys\_user & VARCHAR(50) & 登录账号，唯一 \\
        \hline
        角色 & sys\_user & VARCHAR(20) & ADMIN/STUDENT/TEACHER \\
        \hline
        论文状态 & thesis & VARCHAR(20) & DRAFT/COMPLETED/SUBMITTED/REVIEWED/RETURNED \\
        \hline
        章节内容 & thesis\_section & TEXT & HTML格式的章节正文 \\
        \hline
        模板类型 & template & VARCHAR(20) & GRADUATION/COURSE/PROJECT \\
        \hline
        封面字段配置 & template\_version & JSONB & 定义封面显示哪些字段及其排序 \\
        \hline
        批注偏移量 & annotation & INT & 选中文字在章节HTML中的起始字符位置 \\
        \hline
        分数 & review\_record & INT & 0-100整数，教师评定 \\
        \hline
        等级 & review\_record & VARCHAR(10) & 自动换算：$\geq$90优/≥80良/≥70中/≥60及格 \\
        \hline
    \end{tabular}
\end{table}

\subsection{数据采集}

\begin{itemize}
    \item \textbf{用户数据}：学生自主注册（选择学院、设定账号密码）、教师由管理员后台录入或自主注册；
    \item \textbf{模板数据}：管理员在后台管理界面手动配置封面字段、章节结构和格式参数；
    \item \textbf{论文内容数据}：学生在Web编辑器中实时编辑，通过自动保存（每30秒）和手动保存两种方式持久化到数据库；
    \item \textbf{参考文献数据}：学生通过专用表单逐条录入，支持BibTeX格式批量导入；
    \item \textbf{图片数据}：学生在前端通过文件选择器上传，服务端校验格式和大小后存储；
    \item \textbf{批阅数据}：教师在批阅界面选中论文文字添加批注，填写评语和分数后提交；
    \item \textbf{导出文件}：学生点击导出按钮触发服务端异步生成任务，完成后提供下载。
\end{itemize}
% ============================================
% 过渡：引言/任务概述 → 模板管理功能
% ============================================
% ============================================================
% 2_B_template.tex --- 管理员模板管理功能需求
% 编写人：B  日期：2026-07-15
% 职责模块：学院管理 / 模板管理 / 论文CRUD / 章节存取
% ============================================================

\section{管理员模板管理功能需求}

模板是论文自动生成系统的核心基础，定义了论文的封面样式、字体字号、页边距、页眉页脚、目录格式等全部排版规范参数。管理员通过模板管理功能，维护不同学位类型（本科毕业论文、课程论文、项目论文）、不同学院维度的论文模板，确保所有学生生成的论文文档严格符合学校教务处的排版要求。

\subsection{模板管理功能模块划分}

模板管理功能划分为以下五个子模块：

\begin{enumerate}
    \item \textbf{模板基础管理}：模板元数据的创建、编辑、删除和列表查询；
    \item \textbf{模板版本管理}：模板配置变更时自动生成新版本，保留完整历史记录，支持版本切换与回滚；
    \item \textbf{模板排版参数配置}：封面元素、章节结构、字体格式、页边距、页眉页脚等全部排版参数的在线可视化配置；
    \item \textbf{模板预览与校验}：管理员配置模板后实时预览生成效果，系统自动校验配置合法性；
    \item \textbf{学院管理}：维护学校各学院的基础信息（名称、代码），模板按学院维度进行关联绑定。
\end{enumerate}

\subsection{模板基础管理}

\reqitem{FR-B01}{管理员可创建论文模板，填写模板名称（如\textquotedblleft 计算机科学与技术学院本科毕业论文模板\textquotedblright），选择模板类型（GRADUATION 毕业设计 / COURSE 课程论文 / PROJECT 项目报告），选择适用学院（支持多选，也可设为\textquotedblleft 全部学院\textquotedblright 通用模板）。模板创建后自动生成初始版本（V1.0），初始配置为空，需管理员进一步编辑排版参数。}
\reqitem{FR-B02}{管理员可编辑已有模板的元数据（名称、适用学位类型、适用学院范围）。编辑元数据不触发版本号变更；仅当修改排版配置参数时系统自动递增版本号。}
\reqitem{FR-B03}{管理员可删除不再使用的模板。删除操作前，系统自动检测该模板的历史版本是否被现有论文引用：若存在关联的进行中论文（status 不为 COMPLETED），提示\textquotedblleft 该模板下有 N 篇进行中论文，无法删除\textquotedblright 并阻止删除；若无关联论文或仅有已完成论文，允许删除，同时删除所有关联的 TemplateVersion 记录。}
\reqitem{FR-B04}{管理员进入模板列表页面，支持按模板类型（GRADUATION / COURSE / PROJECT）和学院进行筛选。列表展示字段包括：模板名称、模板类型、适用学院、当前生效版本号、创建时间、最后修改时间。点击模板行进入版本管理详情页。}

\subsection{模板版本管理}

\reqitem{FR-B05}{每次管理员修改模板的排版配置（coverFields / chapterStructure / formatConfig 中任一 JSON 字段）并保存时，系统自动执行版本管理流程：将当前生效版本 is\_current 置为 false；解析当前版本号递增小版本（V1.0\textrightarrow V1.1，V1.9\textrightarrow V2.0）；创建新 TemplateVersion 记录并将 is\_current 置为 true。修改模板元数据（名称、类型）不触发版本递增。}
\reqitem{FR-B06}{管理员可在版本历史页面查看某一模板的全部历史版本列表，展示版本号、创建时间、是否为当前生效版本。支持将任一历史版本设为当前生效版本（回滚操作）：将目标版本的 is\_current 置为 true，原当前版本置为 false。回滚操作不影响已使用该模板的历史论文（论文绑定的 template\_version\_id 保持不变）。}
\reqitem{FR-B07}{管理员选择两个历史版本进行差异对比，系统逐字段比对 coverFields、chapterStructure、formatConfig 三个 JSON 字段的内容差异，前端以对比视图展示变更内容：新增部分以绿色高亮标记，删除部分以红色高亮标记，未变更部分灰色显示。}

\subsection{模板排版参数配置}

\reqitem{FR-B08}{管理员进入模板版本编辑页面，可配置页面参数：纸张大小（默认 A4）、上/下/左/右边距（单位 cm）、行间距倍数（1.5 倍 / 2.0 倍）、段落间距（单位 em）。修改后右侧预览区域实时刷新渲染效果。}
\reqitem{FR-B09}{管理员可配置字体参数：正文字体与字号（默认宋体 12pt）、一级标题字体与字号（默认黑体 16pt）、二级标题字体与字号（默认黑体 14pt）、三级标题字体与字号（默认黑体 12pt）、页眉页脚字体与字号。字体列表从系统可用字体中加载。}
\reqitem{FR-B10}{管理员通过可视化编辑器配置封面元素：校徽图片（上传+预览+位置拖拽）、论文标题（字体/字号/水平位置/垂直位置）、作者姓名格式、指导教师信息格式、学院名称、提交日期格式（支持 yyyy年MM月dd日 等自定义格式）。每个元素独立可拖拽调整位置（x/y 百分比坐标）。}
\reqitem{FR-B11}{管理员配置页眉页脚样式：奇数页页眉内容（可选\textquotedblleft 章标题\textquotedblright / \textquotedblleft 论文题目\textquotedblright / 自定义文本）、偶数页页眉内容、页脚页码格式（如\textquotedblleft 第\{page\}页\textquotedblright 或\textquotedblleft -\{page\}-\textquotedblright）、页码位置（居中/居右）。}
\reqitem{FR-B12}{管理员通过章节树编辑器配置模板的章节结构（chapterStructure）：以树形结构展示论文章节骨架，支持添加/删除章节节点、拖拽调整章节层级和顺序、编辑章节标题。一章节点下可包含多个子节节点，子节节点可继续嵌套。模板内置默认章节结构（摘要、目录、引言、正文各章、结论、致谢、参考文献），管理员可在此基础上自定义调整。}

\subsection{模板预览与校验}

\reqitem{FR-B13}{管理员配置模板排版参数后，可点击\textquotedblleft 在线预览\textquotedblright 按钮。系统使用预设的示例论文内容（包含中英文摘要、目录、多级标题正文、图片、表格、参考文献等代表性元素），按照当前模板配置渲染生成完整的预览 PDF。预览 PDF 在浏览器新标签页中展示，支持缩放和翻页。}
\reqitem{FR-B14}{保存模板配置前，系统自动执行合法性校验，校验项包括：封面必填元素（标题、作者、日期）是否已配置；章节结构是否至少包含\textquotedblleft 正文\textquotedblright 类型的章节节点；字体名称是否在系统支持范围内；页边距设置是否在打印机可打印区域内（通常上下左右均不低于 1.5cm）；JSON 字段格式是否正确。校验不通过时，以红色错误提示框列出所有不合规项及其修复建议，阻止保存。}

\subsection{学院管理}

\reqitem{FR-B15}{管理员可新增学院信息，填写学院名称（如\textquotedblleft 计算机科学与技术学院\textquotedblright）和学院代码（如\textquotedblleft CS\textquotedblright，全局唯一，用于系统内部标识和API路由）。学院代码一经创建后不可修改，但学院名称可编辑。}
\reqitem{FR-B16}{管理员可编辑学院名称，或删除学院。删除操作前，系统自动检测该学院是否存在关联的模板或论文记录：若存在关联模板，提示\textquotedblleft 该学院下存在 N 个关联模板，请先处理关联模板后再删除\textquotedblright；若存在进行中论文，提示\textquotedblleft 该学院下存在进行中的论文，无法删除\textquotedblright。无关联数据时允许删除。}
\reqitem{FR-B17}{管理员在创建模板选择适用学院时，从已维护的学院列表中通过多选下拉框选择。支持\textquotedblleft 全部学院\textquotedblright 快捷选项，选中后该模板适用于所有学院（college\_id 设为 NULL）。}

% ============================================
% 过渡：模板管理 → 参考文献/图片/导出功能
% ============================================
% ============================================
% 第3章：参考文献管理 + 图片上传 + 文档导出功能需求
% 负责人：C
% ============================================

\section{参考文献管理功能}

\subsection{功能描述}

参考文献管理模块为学生提供论文写作过程中参考文献条目的录入、编辑、检索和格式化输出能力。学生可以在系统内独立维护自己的参考文献库，在撰写论文正文时通过引用标注与文献条目关联，最终在导出文档时自动生成符合GB/T 7714-2015国家标准的参考文献列表。

\subsection{功能列表}

\begin{enumerate}
    \item \textbf{参考文献录入}：学生逐条录入参考文献信息，包括作者、标题、文献类型（期刊/书籍/会议/学位论文/网络资源）、出版年份、刊物名称、卷号、期号、页码、出版社、会议名称、学位授予单位、访问链接等字段；
    \item \textbf{文献编辑与删除}：支持对已录入的参考文献信息进行修改和删除操作；
    \item \textbf{文献检索}：支持按作者姓名和文献标题的关键词模糊搜索；
    \item \textbf{标准化格式化}：系统自动将文献信息按GB/T 7714-2015标准格式化为引用字符串，支持五种文献类型的格式化输出；
    \item \textbf{格式化列表预览}：学生可预览全部参考文献的格式化列表效果，确认格式正确再导出。
\end{enumerate}

\subsection{业务流程}

\begin{enumerate}
    \item 学生进入参考文献管理页面，系统展示已录入的文献列表（按年份降序排列）；
    \item 学生点击"新增文献"，弹出录入表单，填写作者、标题、文献类型及其他类型特有字段；
    \item 提交保存后，文献条目出现在列表中，系统同时生成并展示该条目的GB/T 7714格式化文本；
    \item 学生可随时编辑或删除已有文献条目；
    \item 导出论文时，系统自动将所有文献按格式化顺序列表写入文档末尾的"参考文献"章节。
\end{enumerate}

\subsection{数据描述}

参考文献数据以结构化方式存储在数据库中，每条文献包含以下核心信息：

\begin{itemize}
    \item \textbf{基本字段}：作者列表（多作者以逗号分隔）、文献标题、文献类型枚举（J/M/C/D/EB）、出版年份；
    \item \textbf{期刊专用（[J]）}：刊物名称、卷号、期号、页码；
    \item \textbf{书籍专用（[M]）}：出版社名称、出版地；
    \item \textbf{会议专用（[C]）}：会议名称；
    \item \textbf{学位论文专用（[D]）}：学位授予单位；
    \item \textbf{网络资源专用（[EB]）}：访问链接、引用日期。
\end{itemize}

\section{图片上传与管理功能}

\subsection{功能描述}

图片上传管理模块为学生提供论文写作过程中插入图片的功能支持。学生可从本地上传论文插图，系统对上传的图片进行格式校验和大小限制，存储后提供访问URL供前端在编辑器中展示，并在导出文档时将图片嵌入到对应的文档位置。

\subsection{功能列表}

\begin{enumerate}
    \item \textbf{图片上传}：学生通过文件选择器选择本地图片文件，上传至服务器；
    \item \textbf{格式校验}：系统自动校验上传文件的MIME类型，仅允许JPEG、PNG、GIF、WebP格式；
    \item \textbf{大小限制}：单张图片大小限制为10MB，超出限制时返回错误提示；
    \item \textbf{自动重命名}：上传后系统自动以UUID重命名文件，避免文件名冲突；
    \item \textbf{图片预览与访问}：上传成功后返回图片的访问URL，前端可直接嵌入编辑器或展示，支持在浏览器中直接查看图片文件；
    \item \textbf{图片删除}：学生可删除不再需要的已上传图片，同时清理服务器本地文件。
\end{enumerate}

\subsection{业务流程}

\begin{enumerate}
    \item 学生在编辑器中点击"插入图片"按钮，弹出文件选择对话框；
    \item 选择本地图片文件后，前端通过multipart/form-data格式POST请求上传至服务端；
    \item 服务端接收文件，校验MIME类型和文件大小，校验通过后以UUID重命名文件并保存至服务器本地磁盘；
    \item 服务端将图片元信息（原始名称、存储名称、文件路径、大小、类型）写入数据库，返回图片ID和访问URL；
    \item 前端使用返回的URL在编辑器中展示图片，用户保存后该URL嵌入论文章节的HTML内容中；
    \item 导出文档时，服务端根据URL读取本地图片文件，嵌入到DOCX或PDF的对应位置。
\end{enumerate}

\section{文档导出功能}

\subsection{功能描述}

文档导出模块是本系统的核心功能，负责将学生在Web编辑器中完成的论文内容，按照所选模板格式配置，生成标准DOCX（Word）和PDF格式文件供用户下载。导出过程为服务端异步操作，导出完成后提供文件下载链接。

\subsection{功能列表}

\begin{enumerate}
    \item \textbf{DOCX格式导出}：将论文标题、各章节段落、图片、表格、参考文献等内容，结合模板的字体、字号、行距、页边距等格式参数，通过Apache POI库生成标准.docx格式文件；
    \item \textbf{PDF格式导出}：在生成DOCX文件的基础上，通过LibreOffice无头转换模式将DOCX转换为PDF格式，保证排版一致性；
    \item \textbf{导出进度反馈}：导出任务执行期间，系统向前端反馈生成进度（生成中/完成/失败）；
    \item \textbf{文件下载}：导出完成后提供文件下载链接或直接流式传输文件。
\end{enumerate}

\subsection{导出内容范围}

导出的文档应完整包含以下内容：

\begin{enumerate}
    \item \textbf{封面}：论文标题、学生姓名、学院、专业、指导教师、日期等封面字段；
    \item \textbf{原创性声明}：预设的学术诚信声明文本；
    \item \textbf{中英文摘要}：论文的中文和英文摘要内容及关键词；
    \item \textbf{正文章节}：完整的多级章节结构，包含段落文字、图片、表格等素材；
    \item \textbf{参考文献}：按照GB/T 7714-2015标准格式排列的参考文献列表。
\end{enumerate}

\subsection{非功能性需求}

\begin{enumerate}
    \item \textbf{导出质量}：导出的DOCX文件使用宋体、Times New Roman等标准学术字体，中文内容无乱码或方框；
    \item \textbf{响应时间}：单次导出操作应在60秒内完成（论文规模不超过100页、图片10张以内）；
    \item \textbf{并发控制}：同一用户同一时间仅允许执行一次导出操作，防止重复提交；
    \item \textbf{错误处理}：导出失败时需明确告知失败原因（如图片文件缺失、磁盘空间不足等），并提供重试入口；
    \item \textbf{文件存储}：导出的文档文件按日期和用户ID组织存储，定期清理过期文件（保留7天）。
\end{enumerate}
% ============================================
% 过渡：导出功能 → 学生端功能需求
% ============================================
% ============================================================
% 4_D_student.tex — 学生端功能需求 + 交互流程
% 编写人：D  日期：2026-07-15
% ============================================================

\section{学生端功能需求}

\subsection{学生端功能模块划分}

学生端是本系统面向论文撰写者（学生）的核心入口，承担论文创建、编辑、格式化、导出全流程操作。学生端功能划分为以下六大子模块：

\begin{enumerate}
    \item \textbf{账号与个人中心}：注册、登录、个人信息维护、密码修改；
    \item \textbf{论文项目管理}：新建论文、论文列表查看、论文信息编辑、论文删除；
    \item \textbf{章节在线编辑}：富文本章节编写、章节拖拽排序、章节增删、自动保存；
    \item \textbf{内容元素管理}：参考文献录入与格式化引用、图片上传与插入、表格创建与编辑；
    \item \textbf{模板选择与预览}：模板浏览、模板套用、论文实时预览；
    \item \textbf{导出下载}：论文导出为DOCX、PDF格式文件，下载至本地。
\end{enumerate}

\subsection{学生账号与个人中心}

\reqitem{FR-D01}{学生可自主注册账号，填写学号、姓名、邮箱、密码，系统校验学号唯一性及邮箱格式合法性，注册成功后自动跳转登录页。}
\reqitem{FR-D02}{学生使用学号+密码登录系统，后端返回JWT Token，前端存入localStorage，后续请求统一携带Token鉴权；Token过期后自动跳转登录页。}
\reqitem{FR-D03}{学生登录后进入个人中心，可查看和修改姓名、邮箱、手机号等个人信息；修改密码需先输入原密码验证。}
\reqitem{FR-D04}{学生可查看自己创建的全部论文列表，显示论文标题、创建时间、最后修改时间、当前字数统计。}

\subsection{论文项目管理}

\reqitem{FR-D05}{学生点击"新建论文"，进入论文创建向导：（1）输入论文标题；（2）选择所属学院；（3）选择论文模板；（4）确认创建，系统自动按模板生成空白论文骨架（含预设章节结构）。}
\reqitem{FR-D06}{学生在论文列表页可对每篇论文执行：打开编辑、重命名、查看预览、删除操作。删除论文需二次弹窗确认，确认后论文及所有章节数据永久删除。}
\reqitem{FR-D07}{论文列表支持按创建时间、最后修改时间排序，支持按论文标题关键词搜索过滤。}

\subsection{章节在线编辑}

\reqitem{FR-D08}{进入论文编辑器后，左侧显示章节树形导航，右侧为富文本编辑区域。点击章节名称切换编辑目标章节，编辑区同步刷新。}
\reqitem{FR-D09}{编辑器提供完整富文本工具栏：加粗、斜体、下划线、删除线、标题层级（H1-H4）、文字颜色/背景色、段落对齐、有序/无序列表、缩进、撤销/重做。}
\reqitem{FR-D10}{支持在章节正文中通过工具栏按钮或快捷键插入：参考文献引用标记、图片、表格、脚注。引用标记在导出时自动替换为格式化引用编号。}
\reqitem{FR-D11}{学生可在章节树中右键或拖拽调整章节顺序，支持添加同级章节、添加子章节、删除章节、重命名章节。删除章节时若包含子章节，需提示用户一并删除。}
\reqitem{FR-D12}{编辑器每30秒自动保存当前编辑内容至后端，同时监听页面关闭事件，若存在未保存修改则弹窗提示用户。学生也可手动点击"保存"按钮立即保存。}

\subsection{参考文献管理}

\reqitem{FR-D13}{学生可进入论文参考文献管理面板，手动添加参考文献条目，填写：作者、标题、期刊/出版社、年份、卷号/期号、页码、DOI/URL。}
\reqitem{FR-D14}{系统支持常见参考文献类型：期刊论文[J]、会议论文[C]、书籍[M]、学位论文[D]、网络资源[EB/OL]。学生选择类型后，表单字段动态调整。}
\reqitem{FR-D15}{参考文献列表支持编辑、删除、拖拽排序。在正文中通过插入引用标记并选择对应文献，系统按出现顺序自动编号。}
\reqitem{FR-D16}{导出DOCX/PDF时，系统自动在文末生成格式化参考文献列表，格式遵循GB/T 7714-2015国家标准。}

\subsection{模板选择与预览}

\reqitem{FR-D17}{新建论文时或编辑过程中，学生可浏览管理员发布的所有可用模板，按学院筛选。每个模板展示封面缩略图、适用学院、页眉页脚样式说明。}
\reqitem{FR-D18}{学生选择模板后，系统实时将当前论文内容套用新模板样式进行预览，包括字体、行距、页边距、页眉页脚。预览满意后学生确认切换模板。}
\reqitem{FR-D19}{论文编辑过程中，学生可随时点击"预览"按钮，在独立页面查看完整论文排版效果，预览页与实际导出版本视觉效果一致。}

\subsection{导出下载}

\reqitem{FR-D20}{学生点击"导出"，弹出导出选项对话框：选择导出格式（DOCX / PDF），选择导出范围（全部章节 / 指定章节范围），确认后后端生成文件。}
\reqitem{FR-D21}{导出任务提交后，前端显示生成进度（排队中→生成中→完成），完成后自动触发浏览器下载。若生成失败，显示具体错误原因并允许重试。}
\reqitem{FR-D22}{系统保留最近3次导出记录，学生可在导出历史中重新下载之前生成的文件，避免重复生成。}

% ============================================================
\section{学生端核心交互流程}

\subsection{注册登录流程}

\begin{enumerate}
    \item 学生访问系统首页，未登录状态下显示登录/注册入口及系统简介；
    \item 首次使用点击"注册"，填写学号、姓名、邮箱、密码、确认密码；
    \item 前端校验：学号非空且为数字、姓名非空、邮箱格式合法、密码长度≥6位、两次密码一致；
    \item 提交注册请求至 \texttt{POST /api/v1/auth/register}，后端校验学号唯一性，返回注册结果；
    \item 注册成功自动跳转登录页，输入学号+密码提交至 \texttt{POST /api/v1/auth/login}；
    \item 登录成功返回JWT Token，前端存储并跳转至论文列表首页；
    \item 登录失败提示具体错误（账号不存在 / 密码错误 / 账号已禁用）。
\end{enumerate}

\subsection{论文创建与编辑主流程}

\begin{enumerate}
    \item 学生在论文列表页点击"新建论文"，弹出创建向导对话框；
    \item 步骤一：输入论文标题（必填，限200字），选择所属学院（下拉列表）；
    \item 步骤二：系统展示可用模板列表，学生点击模板卡片预览样式；
    \item 步骤三：确认创建，后端创建论文记录，根据所选模板初始化章节结构；
    \item 前端跳转至论文编辑器页面（路由 \texttt{/editor/:paperId}），左侧加载章节树，右侧加载第一个章节内容；
    \item 学生在编辑器中进行富文本编辑，30秒自动保存或手动保存；
    \item 学生点击"预览"查看排版效果，新标签页展示渲染后的HTML；
    \item 学生点击"导出"，选择格式和范围，提交导出任务，等待生成完成下载文件。
\end{enumerate}

\subsection{参考文献插入流程}

\begin{enumerate}
    \item 学生在编辑器工具栏点击"插入引用"按钮，弹出参考文献管理面板；
    \item 面板分为两栏：左侧为已添加的参考文献列表，右侧为添加/编辑表单；
    \item 学生可新增、编辑、删除文献条目，操作实时保存至后端；
    \item 学生选中一条或多条文献，点击"插入引用标记"，系统在编辑器光标位置插入 [1] 占位标记；
    \item 正文编辑过程中引用标记与参考文献列表保持同步：删除文献时提示正文中存在引用，新增文献后引用编号自动重排；
    \item 导出时后端遍历正文引用标记，按出现顺序编号，文末生成标准格式参考文献列表。
\end{enumerate}

\subsection{论文导出流程}

\begin{enumerate}
    \item 学生点击编辑器顶部"导出"按钮，弹出导出配置对话框；
    \item 选择导出格式：DOCX / PDF（单选）；
    \item 选择导出范围：全部章节 / 自定义范围（多选章节）；
    \item 可选勾选：是否包含封面页、是否包含目录、是否包含参考文献列表；
    \item 提交导出请求，后端返回任务ID，前端轮询获取进度；
    \item 状态流转：PENDING → PROCESSING → COMPLETED / FAILED；
    \item 生成完成返回文件下载URL，前端自动触发下载；
    \item 生成失败显示错误阶段及原因，允许重新导出。
\end{enumerate}

% ============================================================
\section{学生端界面原型说明}

\begin{itemize}
    \item \textbf{登录/注册页}：居中卡片式布局，左侧系统Logo及简介，右侧表单区域。登录与注册通过Tab切换，避免页面跳转。
    \item \textbf{论文列表页（首页）}：顶部导航栏含Logo、搜索框、用户头像下拉菜单（个人中心、退出登录）。主体区域为论文卡片网格或列表视图，每张卡片显示论文标题、模板名称、最后修改时间、字数。右上角"新建论文"按钮。
    \item \textbf{论文编辑器页}：三栏布局——左侧可收起的章节树面板（支持拖拽排序），中间主编辑区（富文本编辑器+顶部工具栏），右侧可收起的内容面板（参考文献管理、模板预览切换）。底部状态栏显示自动保存状态和字数统计。
    \item \textbf{个人中心页}：Tab切换个人信息编辑、修改密码、导出历史三个子页面。
\end{itemize}
% ============================================
% 过渡：学生端 → 管理员端与教师端功能需求
% ============================================
% ============================================================
% 5_E_admin_teacher.tex — 管理员端 + 教师端功能需求
% 编写人：E  日期：2026-07-15
% 职责模块：管理员后台管理 / 教师批阅管理 / 账号管理
% ============================================================

\section{管理员端功能需求}

\subsection{管理员端功能模块划分}

管理员端承担系统基础数据的维护和论文模板的配置管理职责，是系统的后台管理核心。管理员端划分为以下几个子模块：

\begin{enumerate}
    \item \textbf{学院信息管理}：对学校各学院信息进行增删改查，学院信息是用户注册和模板绑定的基础；
    \item \textbf{用户账号管理}：查看系统全部用户列表，支持禁用/启用用户账号，查询用户注册信息；
    \item \textbf{模板管理}：创建、编辑、启用/停用论文模板，管理模板版本；
    \item \textbf{系统数据概览}：查看系统运行数据总览（用户数量、论文数量、模板数量等）。
\end{enumerate}

\subsection{学院信息管理}

\reqitem{FR-E01}{管理员可创建新的学院信息，填写学院名称（如"计算机科学与技术学院"）和学院代码（如"CS"，全局唯一）。学院代码一经创建不可修改，学院名称可编辑。}
\reqitem{FR-E02}{管理员可查看全部学院列表，支持按学院名称搜索。列表展示学院名称、学院代码、关联模板数量、关联用户数量。}
\reqitem{FR-E03}{管理员可编辑已有学院名称，或删除无关联数据的学院。删除前系统自动校验该学院是否关联了模板或用户，若存在关联则阻止删除并提示关联数量。}

\subsection{用户账号管理}

\reqitem{FR-E04}{管理员可查看系统全部用户列表，支持按角色（ADMIN/STUDENT/TEACHER）和学院进行筛选。列表展示用户名、真实姓名、角色、所属学院、注册时间、账号状态。}
\reqitem{FR-E05}{管理员可将异常账号（如恶意注册、长期未使用的账号）禁用。禁用后该账号无法登录系统，已登录的Token在Redis中的缓存被清除。支持重新启用已禁用的账号。}
\reqitem{FR-E06}{管理员可重置指定学生的密码为初始密码（123456），重置后该学生下次登录系统时会收到修改密码的提示。}

\subsection{系统数据概览}

\reqitem{FR-E07}{管理员登录后进入系统数据概览仪表盘页面，展示以下统计数据：系统注册用户总数（按角色分类统计）、论文模板总数（按类型分类统计）、在编论文总数（按状态分类统计）、当日导出任务数。数据以卡片和图表形式展示，支持按时间范围筛选。}

% ============================================================
\section{教师端功能需求}

\subsection{教师端功能模块划分}

教师端是面向论文批阅者的操作入口，承载论文审阅、批注添加、评分与退回等核心功能。教师端划分为以下几个子模块：

\begin{enumerate}
    \item \textbf{批阅任务管理}：查看待批阅论文列表、进行中批阅、已完成批阅历史；
    \item \textbf{论文在线批阅}：在线查看论文全文内容，在指定文字位置添加批注，查看已有批注列表；
    \item \textbf{评分与退回}：对论文进行综合评分和等级评定，或退回论文要求学生修改重提；
    \item \textbf{批阅历史查询}：查看历次批阅记录、分数变化轨迹、批注修改对比。
\end{enumerate}

\subsection{批阅任务管理}

\reqitem{FR-E08}{教师登录后进入批阅任务列表页，展示已提交至该教师的全部论文。列表支持按学院筛选、按提交时间排序，展示字段包括：论文标题、学生姓名、所属学院、提交时间、当前状态（待批阅/已批阅/已退回）。}
\reqitem{FR-E09}{教师可点击论文进入批阅详情页。支持在同一页面中顺序浏览多个批阅任务（上一个/下一个），提升批量批阅效率。}

\subsection{论文在线批阅}

\reqitem{FR-E10}{教师进入批阅页面后，左侧展示论文的章节树，右侧展示论文全文HTML渲染内容（只读模式，不可编辑）。教师可通过点击章节树切换查看各章节内容。}
\reqitem{FR-E11}{教师可在论文正文中选中任意文字段落，选中的文字以高亮背景标记，同时弹出批注输入框。教师填写批注内容后确认提交，该批注以卡片形式显示在右侧批注列表中。}
\reqitem{FR-E12}{批注以锚点形式与原文位置绑定，教师点击批注列表中的某条批注，页面自动滚动并高亮对应的原文位置。批注支持编辑和删除操作。}

\subsection{评分与退回}

\reqitem{FR-E13}{教师完成全部批注后，点击"评阅完成"按钮弹出评分对话框。评分要素包括：整体评语（可选填，支持HTML格式）、打分（0-100整数滑块）、等级（系统自动换算）。}
\reqitem{FR-E14}{教师可以选择"通过"（REVIEWED）或"退回"（RETURNED）两种操作。选择退回时必须填写退回原因，且字数不少于10个字符。退回后论文状态变为RETURNED，is\_locked改为FALSE，学生可重新编辑。}
\reqitem{FR-E15}{论文通过后状态变为REVIEWED，学生可查看教师的评语、得分和等级，以及全部批注的详细内容。}}

\subsection{批阅历史查询}

\reqitem{FR-E16}{教师可查看某篇论文的历次批阅记录列表，包括每次批阅的分数、等级、操作类型（通过/退回）、操作时间和对应的退回原因。支持两次批阅的得分变化对比展示。}
\reqitem{FR-E17}{学生可查看历次批阅的全部批注内容和评语，系统在界面上清晰标注哪些批注是本次新增的、哪些是已读的。}
% ============================================
% 过渡：功能需求 → 非功能性需求
% ============================================
% ============================================
% F 负责：非功能性需求（性能 + 安全 + 可维护性）
% ============================================

\section{非功能性需求}

非功能性需求定义了系统在功能正确性之外必须满足的质量属性与约束条件，是保障系统在生产环境中稳定、安全、高效运行的基础。本章从性能、安全性、可维护性三个维度进行详细阐述。

% ===== 3.1 性能需求 =====
\subsection{性能需求}

性能需求关注系统在高并发场景下的响应能力与资源利用效率，是本系统支撑多用户同时编辑、提交与批阅论文的核心质量指标。

\subsubsection{响应时间}

\begin{itemize}
    \item \textbf{页面加载：} 前端页面首屏加载时间不超过 2 秒（含 SPA 资源初始化），在 4G 网络条件下不超过 3 秒。
    \item \textbf{API 响应：} 95\% 的 RESTful API 请求在 500ms 内返回结果；涉及 DOCX/PDF 导出的接口因需调用 LibreOffice 或 Apache POI 进行文件格式转换，允许上限为 10 秒，并采用异步任务队列（如基于 Redis 的消息队列）避免阻塞主线程。
    \item \textbf{数据库查询：} 单表 CRUD 操作在 100ms 内完成；含多表 JOIN 的复杂查询（如章节树与批注联合加载）在 300ms 内完成。
\end{itemize}

\subsubsection{并发能力}

\begin{itemize}
    \item \textbf{并发用户数：} 系统应支持至少 200 个并发用户同时在线编辑论文、提交论文与查看批注。
    \item \textbf{连接池配置：} 数据库连接池（HikariCP）最大连接数设置为 20，Redis 连接池（Lettuce）最大连接数设置为 50，确保在并发峰值时不耗尽数据库与缓存资源。
    \item \textbf{限流策略：} 对登录接口、导出接口等高频/高消耗端点实施令牌桶限流（基于 Redis 计数器），单用户每秒请求数不超过 10 次。
\end{itemize}

\subsubsection{缓存策略}

\begin{itemize}
    \item \textbf{热点数据缓存：} 学院列表、模板结构、参考文献格式模板等读多写少的数据使用 Redis 缓存，减少数据库压力，缓存命中率目标 $\ge 85\%$。
    \item \textbf{会话缓存：} 用户登录后的 JWT Token 黑名单与刷新令牌存储于 Redis，实现无状态服务的快速鉴权与登出。
    \item \textbf{章节草稿缓存：} 学生编辑论文时，前端定时（每 30 秒）将草稿内容通过 Redis 暂存，实现多端切换不丢失编辑进度。
\end{itemize}

\subsubsection{吞吐量}

\begin{itemize}
    \item 系统整体 TPS（每秒事务数）在混合负载（80\% 读 + 20\% 写）下不低于 500。
    \item DOCX/PDF 导出吞吐量不低于 20 份/分钟（单机部署）。
\end{itemize}

% ===== 3.2 安全性需求 =====
\subsection{安全性需求}

安全性需求确保系统能够抵御常见 Web 攻击，保护用户数据与论文内容的机密性、完整性与可用性。

\subsubsection{身份认证}

\begin{itemize}
    \item \textbf{认证方式：} 采用 JWT（JSON Web Token）无状态认证机制。用户使用用户名+密码登录后，服务端签发有效期 24 小时的 Access Token，前端将其存储于 HttpOnly Cookie（非 localStorage）以防范 XSS 攻击窃取令牌。
    \item \textbf{密码策略：} 用户密码使用 BCrypt 算法加盐哈希存储，盐值随机生成，杜绝明文或弱哈希存储。密码最小长度 8 位，须包含字母与数字。
    \item \textbf{会话管理：} 服务端维护 Redis 中的 Token 黑名单，用户主动登出或修改密码后将当前 Token 加入黑名单，实现即时失效。Token 续期采用 Refresh Token 机制，减少重复登录。
\end{itemize}

\subsubsection{访问控制}

\begin{itemize}
    \item \textbf{基于角色的访问控制（RBAC）：} 系统定义三种角色——系统管理员（ADMIN）、教师（TEACHER）、学生（STUDENT）。通过自定义注解 \texttt{@RoleRequired} 与拦截器 \texttt{RoleInterceptor} 对每个 API 端点进行角色校验。
    \item \textbf{数据级权限：} 学生仅可查看与编辑本人的论文；教师仅可批阅所属学院的论文；管理员拥有全局管理权限。所有数据查询均在后端通过当前登录用户 ID 进行过滤，不依赖前端传参。
    \item \textbf{API 防护：} 除认证接口外，所有 \texttt{/api/**} 路径均经过 JWT 过滤器校验，非法请求在网关层即被拦截并返回 HTTP 401。
\end{itemize}

\subsubsection{数据安全}

\begin{itemize}
    \item \textbf{传输安全：} 生产环境强制启用 HTTPS，所有 API 通信经 TLS 1.2+ 加密。开发环境可暂用 HTTP。
    \item \textbf{输入校验：} 所有用户输入（表单、文件上传、URL 参数）在后端进行二次校验，防止 SQL 注入、XSS、CSRF 攻击。Spring Boot 集成 Hibernate Validator 对 DTO 字段进行注解式校验。
    \item \textbf{文件上传安全：} 上传文件限制为图片类型（jpg/png/gif），最大单文件 10MB，通过文件魔数（Magic Number）校验文件真实类型，禁止可执行文件与脚本文件上传。上传文件存储于非 Web 可访问目录，通过专用接口代理访问。
    \item \textbf{SQL 注入防护：} 使用 Spring Data JPA 的参数化查询，杜绝字符串拼接 SQL。MyBatis 或原生 SQL 场景强制使用 \texttt{PreparedStatement}。
\end{itemize}

\subsubsection{日志与审计}

\begin{itemize}
    \item \textbf{操作日志：} 关键操作（登录/登出、论文提交、批阅打分、权限变更）记录至数据库审计表，包含操作人、操作时间、IP 地址、操作内容。
    \item \textbf{敏感信息脱敏：} 日志中不得记录明文密码、Token 等敏感信息。异常堆栈中涉及的敏感参数进行脱敏处理。
\end{itemize}

% ===== 3.3 可维护性需求 =====
\subsection{可维护性需求}

可维护性需求保障系统在交付后能够高效地进行缺陷修复、功能扩展与环境迁移。

\subsubsection{代码规范与结构}

\begin{itemize}
    \item \textbf{分层架构：} 严格遵循 Controller → Service → Repository 三层架构，各层职责清晰：Controller 处理 HTTP 请求与参数校验，Service 编排业务逻辑，Repository 封装数据访问。
    \item \textbf{统一返回格式：} 所有 API 响应使用统一的 \texttt{Result<T>} 包装类，包含状态码（code）、消息（message）、数据（data）三个字段，前端可据此统一处理成功与异常分支。
    \item \textbf{全局异常处理：} 通过 \texttt{@RestControllerAdvice} 的 \texttt{GlobalExceptionHandler} 统一捕获并处理业务异常（\texttt{BusinessException}）与系统异常，向前端返回结构化的错误信息，避免异常堆栈泄露。
    \item \textbf{命名规范：} 遵循阿里巴巴 Java 开发手册规范。类名采用大驼峰（PascalCase），方法名与变量名采用小驼峰（camelCase），数据库表名与字段名采用下划线分隔（snake\_case）。
\end{itemize}

\subsubsection{配置管理}

\begin{itemize}
    \item \textbf{外部化配置：} 数据库连接、Redis 地址、JWT 密钥、文件上传路径等环境相关配置均存放于 \texttt{application.properties}，支持通过环境变量覆盖，实现一次构建、多环境部署。
    \item \textbf{多环境 Profile：} 提供 dev（开发）、test（测试）、prod（生产）三套 Spring Profile，分别对应不同的数据库实例与日志级别。
\end{itemize}

\subsubsection{可测试性}

\begin{itemize}
    \item \textbf{单元测试：} 核心业务逻辑（Service 层）编写 JUnit 5 + Mockito 单元测试，覆盖率目标 $\ge 70\%$。
    \item \textbf{接口测试：} 使用 Spring MockMvc 对 Controller 层进行集成测试，覆盖正常流程与异常边界。
    \item \textbf{Swagger 文档：} 集成 SpringDoc OpenAPI（Swagger 3），自动生成 API 文档页面（\texttt{/doc.html}），方便前后端联调与手工测试。
\end{itemize}

\subsubsection{部署与运维}

\begin{itemize}
    \item \textbf{容器化部署：} 提供 Dockerfile 与 docker-compose.yml，实现应用服务 + PostgreSQL + Redis 的一键编排部署，降低环境差异导致的部署问题。
    \item \textbf{健康检查：} 集成 Spring Boot Actuator，暴露 \texttt{/actuator/health} 端点用于容器编排平台的存活探针与就绪探针。
    \item \textbf{数据库迁移：} 使用 JPA 的 \texttt{ddl-auto=update} 策略在开发阶段自动同步表结构；生产环境切换为 \texttt{validate}，结合 Flyway 或手动 SQL 脚本进行版本化数据库迁移。
\end{itemize}

\subsubsection{文档化}

\begin{itemize}
    \item 所有对外 API 接口须在 Swagger 中标注请求方法、路径、参数说明与返回值示例。
    \item 核心模块（缓存策略、批阅流程、部署步骤）提供单独的技术说明文档，降低新成员接手维护的学习成本。
\end{itemize}
\end{document}
